%% ma: L12-B14-0138
%% lop: 12
%% chuong: 5
%% bai: 14
%% dang: 12.14.42
%% loai: DS
%% muc: [NB, TH, TH, TH]
%% dap_an: "ĐĐĐĐ"
%% nguon: "Ảnh giáo viên gửi 25/11/2026 (tài liệu ôn tập Phương trình mặt phẳng), Phần II Câu 3 (điểm A(1;2;3), mặt phẳng (P): x+2y+2z-2=0)"
%% kien_thuc: "Bài 14 (vectơ pháp tuyến, giao điểm với trục, khoảng cách từ điểm đến mặt phẳng, mặt phẳng song song)"
%% trang_thai: cho-duyet
%% ly_do: "Dạng toán mới đề xuất, chờ giáo viên duyệt"
%% ghi_chu: "Ý b) trong đề gốc dùng lại tên A cho giao điểm của (P) với Oz (khác điểm A(1;2;3) ở đầu đề); giữ nguyên theo đề gốc. Đáp án tự giải (sympy)."
\begin{ex}
Trong không gian $Oxyz$, cho điểm $A(1;2;3)$ và mặt phẳng $(P)\colon x+2y+2z-2=0$.
\choiceTF
{\True $\vec n=(1;2;2)$ là một vectơ pháp tuyến của $(P)$.}
{\True $(P)$ cắt $Oz$ tại điểm $A(0;0;1)$.}
{\True Khoảng cách từ $A$ đến mặt phẳng $(P)$ bằng $3$.}
{\True Mặt phẳng qua $A$ và song song với $(P)$ có phương trình là $x+2y+2z-11=0$.}
\loigiai{a) Đúng, vì $(P)$ có hệ số $(1;2;2)$.
b) Giao điểm với $Oz$ có $x=y=0$: $2z-2=0\Rightarrow z=1$, tức điểm $(0;0;1)$. Đúng.
c) $d(A,(P))=\dfrac{|1+4+6-2|}{\sqrt{1+4+4}}=\dfrac93=3$. Đúng.
d) Mặt phẳng song song $(P)$ có dạng $x+2y+2z+d=0$ ($d\ne-2$); qua $A$: $1+4+6+d=0\Rightarrow d=-11$. Đúng.}
\end{ex}
